\documentclass[../../../include/open-logic-section]{subfiles}

\begin{document}
\olfileid{sth}{ord-arithmetic}{intro}

\olsection{مقدمه}

در \olref[ordinals][]{chap} نظریه‌ای برای اعداد ترتیبی پروراندیم.
در \olref[spine][]{chap} دیدیم که می‌توان اعداد ترتیبی را همچون
ستونی دانست که باقیِ سلسله‌مراتب پیرامون آن ساخته می‌شود. اما این
تنها نقش اعداد ترتیبی نیست. باید حسابِ ترتیبی را نیز انجام دهیم.

پیش‌تر در \olref[ordinals][idea]{sec}، هنگامی که از $\omega$،
$\omega+1$ و $\omega+\omega$ سخن گفتیم، به این موضوع اشاره کردیم.
آن هنگام غیررسمی سخن گفتیم؛ اکنون وقت آن است که مطلب را به‌درستی
شرح دهیم. البته باید یادآوری کنیم که در این فصل فلسفهٔ چندانی در
کار نیست؛ تنها بسط‌های فنی‌اند، همراه با این مشاهدهٔ (تا حدی) جالب
که می‌توانیم یک کار را به دو شیوهٔ متفاوت انجام دهیم.

\end{document}
